\documentclass[12pt,a4paper]{article}

% ----------------------------------------
% Paquetes básicos
% ----------------------------------------
\usepackage[utf8]{inputenc}
\usepackage[T1]{fontenc}
\usepackage[spanish]{babel}
\usepackage{graphicx}
\usepackage{hyperref}
\usepackage{amsmath, amssymb}
\usepackage{verbatim}
\usepackage{listings}
\usepackage{xcolor}
\usepackage{geometry}
\usepackage{setspace}
\usepackage{fancyhdr}
\usepackage{titlesec}
\usepackage{caption}
\usepackage{float}
\usepackage{url}
\usepackage[gen]{eurosym}
\usepackage{mathptmx}

% ----------------------------------------
% Definición de variables del documento
% ----------------------------------------
\newcommand{\DocTitle}{Plantilla de Informe Pericial de Análisis Forense Informático}
\newcommand{\DocAuthor}{Roberto Ripoll Cones}
\newcommand{\DocCenter}{CIPFP MISLATA}
\newcommand{\DocCourse}{CE CETI}
\newcommand{\DocSubject}{ANÁLISIS FORENSE INFORMÁTICO}
\newcommand{\DocYear}{2025--2026}
\newcommand{\DocDate}{\today}

% ----------------------------------------
% Configuración de márgenes
% ----------------------------------------
\geometry{ top=3cm, bottom=3cm, left=3cm, right=3cm }
\def\UrlBreaks{\do\/\do-\do.\do_}

% ----------------------------------------
% Interlineado
% ----------------------------------------
\onehalfspacing

% ----------------------------------------
% Encabezado y pie de página
% ----------------------------------------
\pagestyle{fancy}
\fancyhf{}

\fancyhead[L]{\DocCenter}
\fancyhead[C]{\parbox{0.5\textwidth}{\centering \textbf{\DocTitle}}}
\fancyhead[R]{\DocCourse}

\renewcommand{\headrule}{\vspace{6pt}\hrule\vspace{0pt}}

\setlength{\headheight}{50pt}
\setlength{\headsep}{20pt}

\fancyfoot[L]{\DocDate}
\fancyfoot[C]{\DocAuthor}
\fancyfoot[R]{\thepage}

% ----------------------------------------
% Título
% ----------------------------------------
\title{\DocTitle}
\author{\DocAuthor}
\date{\DocDate}

\begin{document}
\sloppy

% ----------------------------------------
% Portada
% ----------------------------------------
\begin{titlepage}
\centering
{\scshape\LARGE \DocCenter \par}
\vspace{0.5cm}
{\scshape\Large \DocCourse \par}
{\scshape\large \DocSubject \par}
\vspace{2cm}
{\Huge\bfseries \DocTitle \par}
\vspace{2cm}
{\large Autor: \DocAuthor \par}
{\large Años: \DocYear \par}
\vfill
{\large \DocDate \par}
\end{titlepage}

% ----------------------------------------
% Índice general
% ----------------------------------------
\tableofcontents
\newpage

% ----------------------------------------
% Índice de figuras y tablas
% ----------------------------------------
\listoffigures
\newpage

\listoftables
\newpage


% ----------------------------------------
% Glosario
% ----------------------------------------
\section*{Glosario}
\addcontentsline{toc}{section}{Glosario}
Definición de términos técnicos relevantes para el informe.
\newpage

% ----------------------------------------
% Lista de Acrónimos
% ----------------------------------------
\section*{Lista de Acrónimos}
\addcontentsline{toc}{section}{Lista de Acrónimos}
Lista de abreviaturas usadas en el documento.
\newpage

% =======================================================
% 1. Identificación
% =======================================================
\section{Identificación}
\begin{itemize}
    \item \textbf{Perito informático:} \rule{10cm}{0.2pt}
    \item \textbf{Número de colegiado (si aplica):} \rule{7cm}{0.2pt}
    \item \textbf{Entidad / Empresa:} \rule{10cm}{0.2pt}
    \item \textbf{Número de caso o referencia:} \rule{7cm}{0.2pt}
    \item \textbf{Fecha de inicio del análisis:} \rule{6cm}{0.2pt}
\end{itemize}

\newpage

% =======================================================
% 2. Estado legal del análisis
% =======================================================
\section{Estado legal del análisis}
Indicar quién solicita el análisis (empresa, juzgado, particular), ámbito (penal, laboral, administrativo), y límites legales aplicables.

\newpage

% =======================================================
% 3. Declaración de tachas
% =======================================================
\section{Declaración de tachas}
El perito declara que no concurre ninguna causa de tacha o impedimento que pueda comprometer su objetividad.

\newpage

% =======================================================
% 4. Juramento o promesa
% =======================================================
\section{Juramento o promesa}
El perito manifiesta que actuará con objetividad, independencia y veracidad conforme a la ley.

\newpage

% =======================================================
% 5. Resumen ejecutivo
% =======================================================
\section{Resumen ejecutivo}
Breve resumen del caso, hallazgos y conclusiones clave.

\newpage

% =======================================================
% 6. Introducción
% =======================================================
\section{Introducción}

\subsection{Alcance}
Describir el alcance del análisis realizado.

\subsection{Objetivos}
Enumerar los objetivos principales del análisis.

\subsection{Metodología}
Describir metodología siguiendo estándares (ISO 27037, NIST, etc.).

\newpage

% =======================================================
% 7. Detalles del caso
% =======================================================
\section{Detalles del caso}

\subsection{Información del sistema analizado}
\begin{itemize}
    \item Tipo de dispositivo:
    \item Sistema operativo:
    \item Versiones:
    \item Identificadores de hardware:
\end{itemize}

\subsection{Evidencias recolectadas}
Indicar tipo, hash, fecha, herramienta y procedimiento.

\newpage

% =======================================================
% 8. Procedimientos de adquisición
% =======================================================
\section{Procedimientos de adquisición}
Explicar paso a paso cómo se obtuvieron las evidencias, herramientas, parámetros, sellado, verificación de integridad, etc.

\newpage

% =======================================================
% 9. Cronología del incidente
% =======================================================
\section{Cronología del incidente}
Tabla con los eventos relevantes:

\begin{table}[H]
\centering
\begin{tabular}{|c|p{10cm}|}
\hline
Fecha/Hora & Evento relevante \\ \hline
 &  \\ \hline
\end{tabular}
\end{table}

\newpage

% =======================================================
% 10. Análisis forense
% =======================================================
\section{Análisis forense}

\subsection{Herramientas utilizadas}
Listar herramientas y versiones.

\subsection{Procedimientos aplicados}
Descripción detallada del análisis realizado.

\subsection{Resultados obtenidos}
Presentación de los hallazgos principales.

\newpage

% =======================================================
% 11. Limitaciones del análisis
% =======================================================
\section{Limitaciones del análisis}
Descripción clara de cualquier limitación técnica o de alcance que pudiera afectar a las conclusiones.

\newpage

% =======================================================
% 12. Conclusiones y recomendaciones
% =======================================================
\section{Conclusiones y recomendaciones}
Conclusiones basadas estrictamente en los resultados.
Recomendaciones de prevención y mitigación.

\newpage

% =======================================================
% 13. Cadena de custodia
% =======================================================
\section{Cadena de custodia}

\begin{table}[H]
\centering
\begin{tabular}{|c|c|c|c|c|}
\hline
Fecha & Hora & Responsable & Acción realizada & Firma \\ \hline
 &  &  &  &  \\ \hline
\end{tabular}
\end{table}

\newpage

% =======================================================
% 14. Anexos
% =======================================================
\section{Anexos}
\begin{itemize}
    \item Capturas de pantalla
    \item Logs relevantes
    \item Hashes
    \item Listados de comandos
    \item Evidencias complementarias
\end{itemize}

\newpage

% =======================================================
% 15. Declaración de integridad
% =======================================================
\section{Declaración de integridad}
El perito certifica que las evidencias no han sido alteradas durante el proceso de análisis.

\vspace{2cm}

\noindent\textbf{Firma del perito:} \\
\rule{8cm}{0.3pt}

\vspace{1cm}

\noindent\textbf{Firma digital (si aplica):} \\
\rule{8cm}{0.3pt}

\end{document}
